\begin{recipe}
[% 
    preparationtime = {\unit[10]{min} + \unit[2]{h}\unit[20]{min} Kühlzeit},
    portion = {\portion{1}},
    source = {\protect\recipesource{https://www.einfachkochen.de/rezepte/erdbeereis-ohne-eismaschine-einfach-selber-machen}{Einfach Kochen: Erdbeereis ohne Eismaschine}},
    calory = {675kcal | 12g Protein | 132g KH | 10g Fett},
    price = {3,10€},
    created = {10.06.2026},
    %rating = {1},
]
{Erdbeereis ohne Eismaschine}

%\graph{% pictures
%    big=pic/dessert/13_erdbeereis-ohne-eismaschine
%}

\introduction{%
Cremiges Erdbeereis ohne Eismaschine aus nur drei Zutaten. Also quasi Sommer aus dem Mixer, nur mit weniger Aufwand als einmal ernsthaft über eine Eismaschine nachzudenken.
}

\ingredients{
    300 g & Erdbeeren\index{Obst!Erdbeere}\\
    175 g & Naturjoghurt\index{Milchprodukte \& Alternativen!Joghurt}\\
    80 g & Zucker
}
https://www.einfachkochen.de/rezepte/erdbeereis-ohne-eismaschine-einfach-selber-machen
\preparation{%
    \step%
    Erdbeeren waschen, putzen und in grobe Würfel schneiden.
    
    \step%
    Erdbeerwürfel in einen verschließbaren Behälter geben und mindestens 2 Stunden einfrieren.
    
    \step%
    Gefrorene Erdbeeren im Mixer oder mit dem Stabmixer kurz vorpürieren.
    
    \step%
    Joghurt und Zucker zugeben und rasch zu einer cremigen Eismasse mixen.
    
    \step%
    Eis in einen Behälter füllen und mindestens 20 Minuten ins Gefrierfach stellen.
    
    \step%
    Direkt servieren oder vor dem Servieren kurz antauen lassen.
}

\hint{
    Eis ohne Eismaschine wird nach längerem Einfrieren sehr fest. Dann am besten etwa 30 Minuten vorher aus dem Gefrierfach nehmen und im Kühlschrank antauen lassen.
}

\suggestion[Portion]{%
    Das Rezept ergibt etwa 8 Kugeln Eis. Als Dessert kann man es auch gut als 4 kleine Portionen rechnen.
}

\end{recipe}